% Generated draft: requires mathematical and source review.
\documentclass[11pt]{article}
\usepackage[margin=25mm]{geometry}
\usepackage{fontspec}
\setmainfont{Latin Modern Roman}
\usepackage{amsmath,amssymb,mathtools,unicode-math}
\setmathfont{Latin Modern Math}
\usepackage{enumitem,needspace,xurl,hyperref}
\setlength{\parindent}{0pt}
\setlength{\parskip}{5pt}
\setlength{\emergencystretch}{3em}
\newcommand{\N}{\mathbb{N}}
\newcommand{\Z}{\mathbb{Z}}
\newcommand{\Q}{\mathbb{Q}}
\newcommand{\R}{\mathbb{R}}
\newcommand{\C}{\mathbb{C}}
\newcommand{\F}{\mathbb{F}}
\newcommand{\A}{\mathbb{A}}
\newcommand{\PP}{\mathbb{P}}
\newcommand{\E}{\mathbb{E}}
\newcommand{\eps}{\varepsilon}
\newcommand{\dd}{\,\mathrm{d}}
\begin{document}
\section*{UnsolvedMath ID: 800304}
\textbf{Source problem code:} catalogue entry
\subsection{Definitions and mathematical statement}
Let \(S^1=\R/\Z\). Under the usual integer-regularity interpretation, let \(r\in\Z\) with \(r\geq 1\), and let \(\mathrm{Diff}_+^r(S^1)\) denote the group of orientation-preserving \(C^r\) diffeomorphisms of \(S^1\), whose inverses are also \(C^r\). Equip this group with the usual \(C^r\) topology, namely uniform convergence of maps and their derivatives through order \(r\) in local coordinates on the compact circle. Define the commuting-pair space
\[
\mathcal C_r=\bigl\{(f,g)\in\mathrm{Diff}_+^r(S^1)\times\mathrm{Diff}_+^r(S^1):f\circ g=g\circ f\bigr\},
\]
and give it the subspace topology inherited from the product \(C^r\) topology.

The question is whether, for every such \(r\), the space \(\mathcal C_r\) is locally connected. Explicitly, is it true that for every \((f,g)\in\mathcal C_r\) and every relatively open set \(U\subseteq\mathcal C_r\) containing \((f,g)\), there exists a relatively open connected set \(V\subseteq\mathcal C_r\) such that
\[
(f,g)\in V\subseteq U?
\]

The supplied statement writes only \(r\geq1\) and does not explicitly specify the topology or the allowed regularity parameters. The formulation above makes the standard finite integer \(C^r\) interpretation explicit; inclusion of noninteger regularity or \(r=\infty\), and the associated topology, remains unresolved by the supplied record.
\subsection{Short English statement}
For every positive integer differentiability order \(r\), is the space of commuting pairs of orientation-preserving circle diffeomorphisms locally connected in the usual \(C^r\) topology? This makes explicit a conventional interpretation of the source's unspecified regularity range and topology. The supplied registry labels the problem open; that status has not been independently verified.
\subsection{Sources}
\begin{itemize}
\item Local UnsolvedMath registry, permanent ID \(800304\), source code \texttt{catalogue entry}, titled ``Group Actions on 1-Manifolds --- Question 21.'' The record attributes the problem to Navas, with proposed year 2017, and labels it open. Its supplied creation and update date is 2026-10-16.
\item Navas, \emph{Group Actions on 1-Manifolds: A List of Very Concrete Open Questions} (2017), \url{https://arxiv.org/abs/1712.06462}. The supplied provenance identifies the item as ``Question 21 (\#20)'', records access on 2026-10-16, and specifies extraction by \texttt{arxiv\_latex}. These are supplied provenance details; neither the source nor the access claim has been independently verified here.
\end{itemize}
\end{document}
